\documentclass[../main.tex]{subfiles}
\ifSubfilesClassLoaded{\setcounter{chapter}{15}}{}
\begin{document}

\chapter{Evaluasi, Refleksi, dan Integrasi Kompetensi}

Bab akhir ini berisi evaluasi, refleksi, dan integrasi kompetensi sesuai struktur OBE. Bab ini memuat kumpulan asesmen: \textbf{Quiz 1} (CPMK-1), \textbf{Quiz 2} (CPMK-2), \textbf{Quiz 3} (CPMK-3), \textbf{Quiz 4} (CPMK-4), \textbf{Tugas}, \textbf{UTS} (CPMK-1 dan CPMK-2), \textbf{UAS} (CPMK-3 dan CPMK-4), serta \textbf{Project} (aplikasi web sistem informasi akademik sederhana, menguji CPMK-3). Quiz, Tugas, UTS, dan UAS berupa \textbf{pilihan ganda}; Project berupa tugas membuat aplikasi web SIA sederhana. Bab ini juga berisi rubrik penilaian komprehensif berbasis OBE, tinjauan pencapaian kompetensi, dan rekomendasi tindak lanjut pembelajaran.

\begin{subcpmk}
  \item Sub-CPMK 1.1--4.3: Asesmen menyeluruh terhadap seluruh kompetensi mata kuliah melalui Quiz 1--4 (CPMK-1 s.d. CPMK-4), Tugas, UTS (CPMK-1 dan CPMK-2), UAS (CPMK-3 dan CPMK-4), serta Project (aplikasi web SIA, CPMK-3).
\end{subcpmk}

\section{Evaluasi, Refleksi, dan Integrasi Kompetensi}

Bab akhir ini merupakan titik pengukuran pencapaian kompetensi secara menyeluruh. Sesuai struktur OBE, bab ini berisi kumpulan asesmen (Quiz 1--4, Tugas, UTS, UAS, Project), rubrik penilaian komprehensif, tinjauan pencapaian kompetensi, dan rekomendasi studi lanjutan.

\subsection{Quiz 1 --- Menguji CPMK-1}

Quiz 1 menguji penguasaan \textbf{CPMK-1} (Analisis Kebutuhan Pengguna dan Perancangan Antarmuka), meliputi Sub-CPMK 1.1, 1.2, dan 1.3. Format: \textbf{pilihan ganda}. Materi rujukan: Bab 2--4 (pengantar pemrograman web, kebutuhan fungsional dan non-fungsional, wireframe, mockup, aksesibilitas, prinsip UX).

\subsection{Quiz 2 --- Menguji CPMK-2}

Quiz 2 menguji penguasaan \textbf{CPMK-2} (Implementasi Solusi Web Front-End), meliputi Sub-CPMK 2.1--2.4. Format: \textbf{pilihan ganda}. Materi rujukan: Bab 5--11 (HTML, CSS, form, Flexbox/Grid, JavaScript, DOM, event handling, validasi form, library/framework front-end).

\subsection{Quiz 3 --- Menguji CPMK-3}

Quiz 3 menguji penguasaan \textbf{CPMK-3} (Pengembangan Sistem Back-End dan Integrasi Data), meliputi Sub-CPMK 3.1, 3.2, dan 3.3. Format: \textbf{pilihan ganda}. Materi rujukan: Bab 12--15 (PHP dasar, pemrosesan form, OOP, session/cookie, MySQL, CRUD, integrasi API).

\subsection{Quiz 4 --- Menguji CPMK-4}

Quiz 4 menguji penguasaan \textbf{CPMK-4} (Evaluasi, Deployment, dan Administrasi Web), meliputi Sub-CPMK 4.1, 4.2, dan 4.3. Format: \textbf{pilihan ganda}. Materi rujukan: Bab 15 (pengujian, performa, keamanan XSS/CSRF, deployment, administrasi web).

\subsection{Tugas}

Tugas merupakan asesmen formatif per bab (Tugas 1--4 sesuai RPS), mengacu pada materi yang telah dibahas. Format: \textbf{pilihan ganda}. Tugas mendukung pencapaian Sub-CPMK secara bertahap sesuai bab yang relevan.

\subsection{UTS --- Menguji CPMK-1 dan CPMK-2}

Ujian Tengah Semester (UTS) dilaksanakan pada \textbf{Pertemuan 8}. UTS menguji \textbf{CPMK-1} dan \textbf{CPMK-2} secara komprehensif (analisis kebutuhan, wireframe/mockup, HTML, CSS, JavaScript, library/framework). Format: \textbf{pilihan ganda}. Materi rujukan: Bab 2--11.

\subsection{UAS --- Menguji CPMK-3 dan CPMK-4}

Ujian Akhir Semester (UAS) dilaksanakan pada \textbf{Pertemuan 16}. UAS menguji \textbf{CPMK-3} dan \textbf{CPMK-4} secara komprehensif (PHP, MySQL, API, pengujian, keamanan, deployment). Format: \textbf{pilihan ganda}. Materi rujukan: Bab 12--15.

\subsection{Project --- Menguji CPMK-3}

Project menguji penguasaan \textbf{CPMK-3}. Bentuk: \textbf{tugas membuat aplikasi web sistem informasi akademik (SIA) sederhana}. Aplikasi harus mencakup: back-end PHP, koneksi MySQL, operasi CRUD, session, integrasi API front-end--back-end. Front-end dapat memanfaatkan materi HTML/CSS/JavaScript dari bab sebelumnya. Soal lengkap dan rubrik project tercantum di Lampiran.

\subsection{Rubrik Penilaian Komprehensif Berbasis OBE}

Penilaian mengikuti prinsip OBE: setiap komponen asesmen dipetakan ke CPMK/Sub-CPMK. Quiz, Tugas, UTS, dan UAS (pilihan ganda) dinilai berdasarkan jumlah jawaban benar. Project dinilai berdasarkan rubrik yang terhubung ke Sub-CPMK 3.1, 3.2, 3.3. \textbf{Bobot (sesuai RPS):} Tugas/Praktikum 35\%, Kuis (Quiz 1--4) 10\%, UTS 20\%, Proyek 20\%, UAS 15\%. Rubrik terperinci tercantum di \textbf{Lampiran A} buku ajar.

\subsection{Tinjauan Pencapaian Kompetensi}

Setelah menyelesaikan seluruh bab, mahasiswa diharapkan mampu:

\begin{itemize}
  \item \textbf{CPMK-1}: Mengidentifikasi kebutuhan, membuat wireframe/mockup, menerapkan aksesibilitas dan UX.
  \item \textbf{CPMK-2}: Membangun antarmuka web responsif dan dinamis dengan HTML5, CSS3, JavaScript, dan library/framework.
  \item \textbf{CPMK-3}: Mengembangkan back-end PHP, mengelola MySQL, mengintegrasikan API (diuji oleh Quiz 3, UAS, dan Project).
  \item \textbf{CPMK-4}: Melakukan pengujian, menerapkan keamanan, melakukan deployment dan administrasi web.
\end{itemize}

Tinjau checklist di setiap bab; identifikasi area yang masih perlu dikembangkan; buat rencana belajar lanjutan sebelum UAS dan Project.

\subsection{Rekomendasi Tindak Lanjut Pembelajaran}

Kelulusan mata kuliah Pemrograman Web menandai awal perjalanan --- teknologi web terus berkembang. Rekomendasi studi lanjutan:

\begin{itemize}
  \item \textbf{Framework back-end}: Laravel, Symfony, atau CodeIgniter untuk pengembangan web PHP tingkat lanjut.
  \item \textbf{Framework front-end}: Vue.js, React, atau Angular untuk aplikasi single-page (SPA).
  \item \textbf{DevOps dan CI/CD}: Docker, GitHub Actions, atau Jenkins untuk otomasi deployment.
  \item \textbf{Keamanan web lanjut}: OWASP Top 10 mendalam, penetration testing dasar, keamanan API.
  \item \textbf{Performa dan skala}: caching (Redis, Memcached), load balancing, optimasi database.
  \item \textbf{Referensi standar}: MDN Web Docs, OWASP, web.dev, PHP The Right Way.
\end{itemize}

Teknologi web terus bergerak; seorang developer harus terus belajar dan mengikuti praktik terbaik (best practices).


\begin{aktivitas}
  \item Pelajari Quiz 1--4: review materi Bab 2--15 sesuai CPMK yang diuji.
  \item Review checklist di setiap bab dan tandai indikator yang sudah dicapai.
  \item Identifikasi area kompetensi yang masih perlu dikembangkan sebelum UAS.
  \item Pelajari rubrik penilaian di Lampiran A dan pastikan Project SIA memenuhi kriteria CPMK-3.
  \item Selesaikan Project: buat aplikasi web sistem informasi akademik sederhana (back-end PHP, MySQL, CRUD, integrasi API).
  \item Susun rencana belajar lanjutan berdasarkan rekomendasi di bab ini.
\end{aktivitas}

\begin{latihan}
  \item Sebutkan komponen asesmen (Quiz 1--4, Tugas, UTS, UAS, Project) dan CPMK yang diuji masing-masing.
  \item Jelaskan keterkaitan UTS dengan CPMK-1 dan CPMK-2.
  \item Jelaskan keterkaitan UAS dengan CPMK-3 dan CPMK-4.
  \item Sebutkan format Quiz, Tugas, UTS, dan UAS (pilihan ganda) serta format Project (aplikasi web SIA).
  \item Buat matriks refleksi: tulis CPMK yang paling dikuasai dan yang perlu ditingkatkan.
  \item Sebutkan tiga rekomendasi studi lanjutan setelah menyelesaikan mata kuliah ini.
\end{latihan}

\begin{asesmen}
Bab ini merupakan dasar asesmen akhir. \textbf{Quiz 1} menguji CPMK-1; \textbf{Quiz 2} menguji CPMK-2; \textbf{Quiz 3} menguji CPMK-3; \textbf{Quiz 4} menguji CPMK-4. Quiz, \textbf{Tugas}, \textbf{UTS}, dan \textbf{UAS} berupa \textbf{pilihan ganda}. \textbf{UTS} (Pertemuan 8) menguji CPMK-1 dan CPMK-2. \textbf{UAS} (Pertemuan 16) menguji CPMK-3 dan CPMK-4. \textbf{Project} menguji CPMK-3: tugas membuat aplikasi web \textbf{sistem informasi akademik (SIA) sederhana} (back-end PHP, MySQL, CRUD, integrasi API).

\textbf{Bobot Penilaian (sesuai RPS):} Tugas/Praktikum 35\%; Kuis (Quiz 1--4) 10\%; UTS 20\%; Proyek 20\%; UAS 15\%. Total 100\%. Rubrik penilaian komprehensif tercantum di \textbf{Lampiran A}.
\end{asesmen}

\begin{checklist}
  \item Saya memahami komponen asesmen: Quiz 1--4 (CPMK-1 s.d. CPMK-4), Tugas, UTS, UAS, Project.
  \item Saya mengetahui format: Quiz, Tugas, UTS, UAS = pilihan ganda; Project = aplikasi web SIA sederhana.
  \item Saya mengetahui UTS menguji CPMK-1 dan CPMK-2; UAS menguji CPMK-3 dan CPMK-4; Project menguji CPMK-3.
  \item Saya dapat mereview checklist tiap bab dan mengidentifikasi area yang perlu ditingkatkan.
  \item Saya mengenal rubrik penilaian komprehensif berbasis OBE di Lampiran A.
  \item Saya mampu merefleksikan pencapaian kompetensi per CPMK.
  \item Saya siap menghadapi UAS dan menyelesaikan Project SIA.
\end{checklist}

\begin{rangkuman}
Bab ini merupakan Bab Akhir OBE yang berisi evaluasi, refleksi, dan integrasi kompetensi. Bab ini memuat kumpulan asesmen: Quiz 1--4 (masing-masing menguji CPMK-1 s.d. CPMK-4), Tugas, UTS (CPMK-1 dan CPMK-2), UAS (CPMK-3 dan CPMK-4), serta Project (aplikasi web SIA sederhana, menguji CPMK-3). Quiz, Tugas, UTS, dan UAS berupa pilihan ganda.

\textbf{Poin Kunci:}
\begin{itemize}
  \item Quiz 1--4 menguji CPMK-1 s.d. CPMK-4; UTS menguji CPMK-1 dan CPMK-2; UAS menguji CPMK-3 dan CPMK-4.
  \item Project: tugas membuat aplikasi web sistem informasi akademik sederhana, menguji CPMK-3.
  \item Rubrik penilaian komprehensif berbasis OBE tercantum di Lampiran A.
  \item Tinjauan pencapaian: review checklist tiap bab, identifikasi area yang perlu ditingkatkan.
  \item Rekomendasi studi lanjutan: framework, DevOps, keamanan, performa, referensi standar.
\end{itemize}
\end{rangkuman}

\ifSubfilesClassLoaded{
  \renewcommand{\bibname}{Daftar Pustaka}
  \bibliographystyle{plain}
  \bibliography{../references}
}{}
\end{document}

